\section{Practical Implications and Deployment Considerations}
\label{sec:praxis}

\subsection{Deployment Controls}
\label{subsec:deployment-controls}

External content should retain information about its origin and trust status throughout
retrieval, summarisation, memory, and tool use. Prompt delimiters or labels alone
remain model-mediated. They become externally enforceable only when the
application preserves them and an independent policy component uses them when
authorising information flows or actions
\parencites{hines2024spotlighting}{wu2024systemlevel}.

Tools and credentials should follow the principle of least privilege. Read
permissions should be separated from write, send, or delete permissions, and
credentials should be short-lived and limited to the current task. An
authorisation component outside the model should validate the requested
operation, recipient, destination, and security-relevant parameters before
execution
\parencites{debenedetti2024agentdojo}
             {debenedetti2025defeating}
             {owasp2025llm}.

Sandboxing can further restrict file, network, code-execution, and secret access
when model output is malicious. It does not prevent prompt injection, but it can
contain the resulting action~\parencites{debenedetti2024agentdojo}{debenedetti2025defeating}{owasp2025llm}. High-impact operations, such as sending messages,
modifying records, or disclosing protected data, should additionally require an
independent policy decision or explicit human approval.  

These controls depend on correct source labels, sufficiently complete
policies, and tool interfaces that cannot bypass the authorisation component.
They reduce the consequences of model failure but do not guarantee that every
permitted action is desirable.

\subsection{Testing, Monitoring, and Incident Response}
\label{subsec:testing-monitoring-response}

Pre-deployment and regression testing should reflect the application's actual
trust boundaries and permissions. Tests should include direct and indirect
payloads, tool interactions, varied attacker objectives, and attacks adapted to
the deployed defence. White-box audits are useful for examining the limits of
open-weight models, but should be distinguished from realistic query-only
attack paths
\parencites{debenedetti2024agentdojo}
             {jia2025critical}
             {yang2025checkpointgcg}
             {pandya2025astra}.

Protected logs should record relevant input origin, retrieved content,
model outputs, proposed tool calls, policy decisions, approvals, credential
scopes, and external destinations. Because prompts and outputs may contain
sensitive information, these records require access control, data minimisation,
retention limits, and integrity protection
\parencite{owasp2025llm}.

An incident plan should support containment of the affected workflow,
revocation of active or exposed credentials, quarantine of malicious content,
and investigation of affected data and actions. The incident should also be
added to future regression tests. Content stored in retrieval systems or memory
may need to be removed before the workflow resumes.

Monitoring and incident response do not prevent the initial model failure and
cannot reverse information that has already been disclosed. Their purpose is to
detect malicious behaviour, contain affected workflows, support investigation,
and reduce continued exposure
\parencites{owasp2025llm}{debenedetti2024agentdojo}.

\subsection{Trade-Offs and Outlook}
\label{subsec:deployment-tradeoffs}

System-level controls introduce operational costs. Policy gates and sandboxes
require integration and maintenance, approval steps add latency and reduce
automation, and monitoring creates storage, privacy, and access-control
obligations. StruQ and SecAlign also require model adaptation and application
support while remaining probabilistic.

Deployment decisions must therefore balance task utility, engineering cost,
policy coverage, and the consequences of both model and enforcement failure.
Where external controls cannot reduce the residual risk to an acceptable level,
high-impact actions should remain subject to independent approval or should not
be automated. Defence in depth combines model-level mitigation with external
prevention, containment, detection, and recovery
\parencites{owasp2025llm}
             {wu2024systemlevel}
             {debenedetti2025defeating}.